%%%%%%%%%%%%%%%%%%%%%%%%%%%%%%%%%
%
%       EXERCÍCIO
%
%%%%%%%%%%%%%%%%%%%%%%%%%%%%%%%%%

\ifdefstring{\atividade}{prova}{%
    \ifdefstring{\modo}{objetivo}{%
        \renewcommand{\valorquestao}{\ValorQObj\ ponto}
    }{%
        \renewcommand{\valorquestao}{\ValorQDisc\ pontos}
    }%
}%

\begin{exercicioBanco}[\valorquestao]
Seja \(M(a)\) o conjunto dos múltiplos de \(a\) e \(D(a)\) o conjunto dos divisores de \(a\), com \(a \in \bbN = \{1, 2, 3, \ldots\}\). Liste os elementos do conjunto \(M(3) \cap D(30)\).

% Define as alternativas
\newcommand{\alternativas}{%
    \begin{center}
        \begin{tabularx}{\textwidth}{XXX}
            (a) \(\{6,15,30\}\). &
            (b) \(\{1,2,3,5,6,10,15,30\}\). &
            (c) \(\{3,6,15,30\}\). \\[5pt]
            (d) \(\{3,15\}\). &
            (e) \(\{3,6,9,12,15,18,21,24,27,30\}\).
        \end{tabularx}
    \end{center}
}

% Define a resposta correta
\newcommand{\resposta}{C}

% Lógica condicional para exibição
\ifdefstring{\atividade}{lista}{%
    \alternativas
    \vspace{0.5em}
    
    \noindent\textbf{Resposta:} letra \textbf{\resposta}.
}{%
    \ifdefstring{\modo}{objetivo}{%
        \alternativas
    }{%
        % modo = discursiva → não mostra alternativas
    }
}
\end{exercicioBanco}

